\itemhead{复习题二 · 第15题（教材p.63）}{XJ-B1-C02-RV-P063-Y0289}
\label{XJ-B1-C02-RV-P063-Y0289}
{\footnotesize 湘教版教材p.63 / PDF第70页 / 复习题二 / 第15题}
\field{原题}{如图，码头 $O$ 南偏东 $60^\circ$ 方向 $400\,\mathrm{km}$ 处（$A$）有一个台风中心。台风正以 $40\,\mathrm{km/h}$ 的速度向正北方向移动，且距台风中心 $350\,\mathrm{km}$ 以内都会受台风影响。问从现在起多长时间后，码头 $O$ 将受台风影响？码头受台风影响的时间大约多久？
\bankfig{C02-RV063}{0.39}}
\field{参考答案}{令经过时间为 $t\ge0$ 小时，距离平方 $=(200\sqrt3)^2+(40t-200)^2$。不超过350千米等价于 $(40t-200)^2\le2500$，故 $t\in[3.75,6.25]$。3.75小时后开始影响，持续2.5小时。}
{\small \field{关联知识}{I04 一元二次不等式与参数分类（教材帮p.64、65、66 / PDF73、74、75）；F03 函数求值、解析式与函数建模（教材帮p.80、87、88、89、90、131 / PDF89、96、97、98、99、140）}
\field{方法与易错点}{化标准式、查系数、求根或判别式、按图定号，最后与实际范围取交集。 易错：套“内小外大”却忘开口；丢重根或降次；漏整数产量限制。}
\field{答案依据}{答案册相关题组：PDF第15、16页（册内页14、15）；按当前题库原题逐项重算，不按旧版题号机械对套。}}
